\section*{Bab \ref{ch:b3:modules} --- Modul atas Daerah Ideal
Utama}

\begin{solution}{exo:b3:modules:1}
(a) Jika $\Q = \Z q_1 + \dots + \Z q_k$, misalkan $d$ sebuah penyebut
persekutuan dari $q_i$: setiap kombinasinya terletak di $\frac1d\Z$,
padahal $\frac1{2d} \notin \frac1d\Z$. Kontradiksi.

(b) \omterm{def:b3:modules:free}{Bebas} \omterm{def:b3:modules:torsion}{torsi}: $nq = 0$ dengan $n \neq 0$ memaksa $q = 0$ di
$\Q$. Tidak \omterm{def:b3:modules:free}{bebas}: dua bilangan rasional tak nol $\frac ab, \frac cd$ mana pun
memenuhi relasi tak trivial $(bc)\frac ab - (ad)\frac cd =
0$, sehingga basisnya beranggota paling banyak satu unsur; sedangkan $\Q \cong \Z$ akan membuat
$\Q = \Z q$ siklik, padahal $\frac q2 \notin \Z q$. (Dan $\Q \neq
0$.)

(c) \cref{thm:b3:modules:structure} mengandaikan \emph{pembangunan
berhingga}, dan (a) menyangkalnya: jadi tak ada kontradiksi --- sebaliknya, $\Q$
menunjukkan bahwa hipotesis itu diperlukan pada pernyataan
``\omterm{def:b3:modules:free}{bebas} \omterm{def:b3:modules:torsion}{torsi} $\Rightarrow$ \omterm{def:b3:modules:free}{bebas}''.
\end{solution}

\begin{solution}{exo:b3:modules:2}
$360 = 2^3\cdot3^2\cdot5$. Partisinya: untuk $3$: $(3), (2,1),
(1,1,1)$; untuk $2$: $(2), (1,1)$; untuk $1$: $(1)$. Karena itu ada $3 \times 2
\times 1 = 6$ grup. \omterm{thm:b3:modules:structure}{Pembagi elementer} $\to$ \omterm{thm:b3:modules:smith}{faktor
invarian}:
\begin{center}
\begin{tabular}{l l}
$\Z/8 \times \Z/9 \times \Z/5$ & $\Z/360$\\
$\Z/8 \times \Z/3 \times \Z/3 \times \Z/5$ & $\Z/3 \times
\Z/120$\\
$\Z/4 \times \Z/2 \times \Z/9 \times \Z/5$ & $\Z/2 \times
\Z/180$\\
$\Z/4 \times \Z/2 \times \Z/3 \times \Z/3 \times \Z/5$ & $\Z/6
\times \Z/60$\\
$(\Z/2)^3 \times \Z/9 \times \Z/5$ & $\Z/2 \times \Z/2 \times
\Z/90$\\
$(\Z/2)^3 \times \Z/3 \times \Z/3 \times \Z/5$ & $\Z/2 \times
\Z/6 \times \Z/30$
\end{tabular}
\end{center}
(Untuk beralih ke \omterm{thm:b3:modules:smith}{faktor invarian}: $d_s$ yang terbesar mengumpulkan
pangkat prima tertinggi dari setiap \omterm{def:b3:rings:divisibility}{unsur prima}, dan seterusnya ke bawah.) Yang berorde
$p^5$: sebanyak partisi dari $5$, yakni $7$.
\end{solution}

\begin{solution}{exo:b3:modules:3}
$B$: $D_1 = \gcd(2,4,6,8) = 2$; $D_2 = \abs{\det B} = \abs{16 -
24} = 8$. \omterm{thm:b3:modules:smith}{Faktor invariannya} $d_1 = 2$, $d_2 = 8/2 = 4$: jadi bentuk
Smithnya $\operatorname{diag}(2, 4)$, dan $\Z^2/B\Z^2 \cong \Z/2\Z
\times \Z/4\Z$.

$C$: diagonal tetapi \emph{bukan} bentuk Smith (sebab $2 \nmid 3$). $D_1 =
\gcd(2,3,12) = 1$; $D_2 = \gcd(2\cdot3,\, 2\cdot12,\, 3\cdot12) =
\gcd(6, 24, 36) = 6$; $D_3 = 72$. Jadi $d = (1, 6, 12)$ dan
$\Z^3/C\Z^3 \cong \Z/6\Z\times\Z/12\Z$ --- sejalan dengan
sisa Tiongkok: $\Z/2\times\Z/3\times\Z/12 \cong \Z/6\times\Z/12$.
\end{solution}

\begin{solution}{exo:b3:modules:4}
Tulis $B = Q\,\operatorname{diag}(d_1, \dots, d_n)\,P$ dengan $Q, P
\in GL_n(\Z)$ (\cref{thm:b3:modules:smith}; tak ada $d_i$ yang nol sebab
$\det B \neq 0$). Maka $\Z^n/B\Z^n \cong \Z^n/
\operatorname{diag}(d)\Z^n = \prod_i \Z/d_i\Z$ (sebab
isomorfisma komposisi $x \mapsto Q^{-1}x$ dari $\Z^n$ memetakan $B\Z^n$ pada
$\operatorname{diag}(d)P\Z^n = \operatorname{diag}(d)\Z^n$).
Kardinalitasnya $\prod \abs{d_i} = \abs{\det
\operatorname{diag}(d)} = \abs{\det B}$, sebab $\det Q, \det P = \pm
1$. Untuk $L = \Z(2,0) + \Z(1,3)$: $B = \bigl(\begin{smallmatrix}2
& 1\\ 0 & 3\end{smallmatrix}\bigr)$, $D_1 = 1$, $D_2 = 6$: jadi $\Z^2/L
\cong \Z/6\Z$, berkardinalitas $\abs{\det B} = 6$.
\end{solution}

\begin{solution}{exo:b3:modules:5}
(a) Submodul: sudah dikerjakan pada \cref{def:b3:modules:torsion}. Jika $a(x +
T(M)) = 0$ di $M/T(M)$ dengan $a \neq 0$, maka $ax \in T(M)$: jadi $bax
= 0$ untuk suatu $b \neq 0$, dan $ba \neq 0$ (di daerah integral), sehingga $x \in
T(M)$: kelasnya nol. Jadi $M/T(M)$ \omterm{def:b3:modules:free}{bebas} \omterm{def:b3:modules:torsion}{torsi}.

(b) \omterm{def:b3:rings:ideal}{Ideal} $(X, Y) \subseteq K[X,Y]$ \omterm{def:b3:modules:free}{bebas} \omterm{def:b3:modules:torsion}{torsi} (sebagai submodul dari
daerah integral $K[X,Y]$ yang bekerja pada dirinya sendiri). Andaikan ia \omterm{def:b3:modules:free}{bebas}; dua
unsur $P, Q$ mana pun memenuhi $Q\cdot P - P \cdot Q = 0$, sebuah relasi tak
trivial bila $P \ne Q$ tak nol, sehingga basisnya beranggota satu unsur:
jadi $(X, Y) = (P)$ berupa \omterm{def:b3:rings:ideal}{ideal} utama --- bertentangan dengan
\cref{exo:b3:rings:6}(a). \omterm{def:b3:modules:free}{Bebas} \omterm{def:b3:modules:torsion}{torsi} dan \omterm{def:b3:modules:free}{dibangun secara berhingga}
(sebab $X, Y$ membangunnya), namun tidak \omterm{def:b3:modules:free}{bebas}: atas $K[X,Y]$ yang bukan \omterm{def:b3:rings:pidufd}{DIU},
teorema strukturnya gagal.
\end{solution}

\begin{solution}{exo:b3:modules:6}
(a) Jika $\Z = 2\Z \oplus C$, proyeksi $\Z \to \Z/2\Z$
terbatas menjadi isomorfisma $C \cong \Z/2\Z$: jadi $C$ akan menjadi
subgrup $\Z$ yang unsur tak nolnya $x$ memenuhi $2x \in C
\cap 2\Z = 0$. Padahal $\Z$ \omterm{def:b3:modules:free}{bebas} \omterm{def:b3:modules:torsion}{torsi}: sehingga $C = 0$, dan itu memaksa $\Z =
2\Z$ --- salah.

(b) \omterm{def:b3:modules:module}{Modul} $A^n/M$ \omterm{def:b3:modules:free}{dibangun secara berhingga} dan \omterm{def:b3:modules:free}{bebas} \omterm{def:b3:modules:torsion}{torsi}, sehingga \omterm{def:b3:modules:free}{bebas}
(\cref{thm:b3:modules:structure}): $A^n/M \cong A^r$ dengan basis
$f_1, \dots, f_r$. Pilihlah prapeta $y_i \in A^n$ dari $f_i$
lalu tetapkan $F = Ay_1 + \dots + Ay_r$. Setiap $x \in A^n$ mempunyai $\pi(x)
= \sum a_if_i$, sehingga $x - \sum a_iy_i \in M$: jadi $A^n = M + F$. Jika
$\sum a_iy_i \in M$, menerapkan $\pi$ memberikan $\sum a_if_i = 0$,
sehingga semua $a_i = 0$ (karena basis): $M \cap F = 0$. Jadi $A^n = M \oplus
F$.
\end{solution}

\begin{solution}{exo:b3:modules:7}
Reduksi pada \cref{exo:b3:modules:3} bersifat efektif: dengan
\[
L = \begin{pmatrix} 1 & 0\\ -3 & 1\end{pmatrix},
\qquad
R = \begin{pmatrix} 1 & 2\\ 0 & -1 \end{pmatrix},
\qquad
LBR = \begin{pmatrix} 2 & 0\\ 0 & 4\end{pmatrix}
\]
(operasi baris $R_2 \leftarrow R_2 - 3R_1$, lalu operasi kolom
$C_2 \leftarrow C_2 - 2C_1$ dan $C_2 \leftarrow -C_2$). Dengan menetapkan
$y = R^{-1}x$ (sebuah bijeksi $(\Z/20\Z)^2$, karena $R$
terbalikkan atas $\Z$), sistem $Bx \equiv b \pmod{20}$ menjadi
ekuivalen dengan
\[
2y_1 \equiv b_1, \qquad 4y_2 \equiv -3b_1 + b_2 \pmod{20}.
\]
Kekongruenan $ky \equiv c \pmod{20}$ \omterm{def:b3:groups:derived}{terselesaikan} jika dan hanya jika
$\gcd(k, 20) \mid c$: jadi penyelesaian ada jika dan hanya jika $2 \mid b_1$ dan $4
\mid b_2 - 3b_1$, yakni \emph{$b_1$ genap dan $b_2 \equiv 3b_1
\pmod 4$}. Bila \omterm{def:b3:groups:derived}{terselesaikan}, ada $2 \times 4 = 8$ penyelesaian
modulo $20$.
\end{solution}

\begin{solution}{exo:b3:modules:8}
(a) Endomorfisma nilpoten $u$ mempunyai $\mu_u = X^k$; \omterm{thm:b3:modules:structure}{pembagi elementernya}
adalah $X^{k_1} \geq \dots$, dengan satu blok Jordan $J_{k_i}(0)$ per bagian
dari sebuah partisi $4$:
\begin{center}
\begin{tabular}{l l l}
partisi & \omterm{thm:b3:modules:jordan}{bentuk Jordan} & \omterm{thm:b3:modules:smith}{faktor invarian}\\
$(4)$ & $J_4$ & $X^4$\\
$(3,1)$ & $J_3 \oplus J_1$ & $X,\ X^3$\\
$(2,2)$ & $J_2 \oplus J_2$ & $X^2,\ X^2$\\
$(2,1,1)$ & $J_2 \oplus J_1 \oplus J_1$ & $X,\ X,\ X^2$\\
$(1,1,1,1)$ & $0$ & $X, X, X, X$
\end{tabular}
\end{center}
(b) Ambil $u = J_2\oplus J_2$ dan $v = J_2 \oplus J_1 \oplus J_1$:
keduanya mempunyai $\chi = X^4$, $\mu = X^2$, tetapi \omterm{thm:b3:modules:smith}{faktor invariannya}
berbeda --- jadi tidak serupa (\cref{thm:b3:modules:frobenius}); kita
juga dapat membandingkan ranknya: $\operatorname{rk} u = 2 \neq 1 =
\operatorname{rk} v$. Untuk $n \leq 3$: $\chi$ dan $\mu$ menentukan,
untuk setiap nilai eigen $\lambda$ (atas sebuah lapangan pemecah), ukuran
total $m_\lambda \leq 3$ dari blok-$\lambda$ beserta blok
terbesarnya $r_\lambda$; dan partisi $m \leq 3$ ditentukan oleh
bagian terbesarnya ($m = 3, r = 2$ memaksa $(2,1)$, dan seterusnya). Jadi
\omterm{thm:b3:modules:structure}{pembagi elementernya} berimpit, dan
\cref{cor:b3:modules:descent} menurunkan keserupaannya ke lapangan
dasarnya.
\end{solution}

\begin{solution}{exo:b3:modules:9}
(i)$\Rightarrow$(ii): jika $V = K[u]x$, maka $V \cong
K[X]/\operatorname{Ann}(x)$, dan $\operatorname{Ann}(x) = (\mu_u)$
(sebab sebuah polinomial mematikan $x$ jika dan hanya jika ia mematikan seluruh $V = K[u]x$, karena
$P(u)Q(u)x = Q(u)P(u)x$). Jadi $n = \dim V = \deg \mu_u$; dan karena
$\mu_u \mid \chi_u$ serta $\deg\chi_u = n$, kemonikan memberikan $\mu_u =
\chi_u$.

(ii)$\Rightarrow$(iii): $\deg\chi_u = \sum_i \deg P_i$ dan
$\mu_u = P_s$ (\cref{thm:b3:modules:frobenius}); kesamaan
derajatnya memaksa $s = 1$.

(iii)$\Rightarrow$(i): $V \cong K[X]/(P_1)$ bersifat siklik, dibangun
oleh prapeta $\bar 1$.

\omterm{def:b3:modules:companion}{Matriks pendamping} merupakan kasus $V = K[X]/(P)$ itu sendiri: $x = \bar
1$, yakni $e_1$, bersifat siklik. Untuk matriks diagonal
$\operatorname{diag}(\lambda_1, \dots, \lambda_n)$: $\chi = \prod
(X - \lambda_i)$, $\mu = \prod_{\lambda \text{ berbeda}} (X -
\lambda)$; keduanya bersesuaian jika dan hanya jika $\lambda_i$ berbeda berpasangan:
jadi matriks diagonal mempunyai vektor siklik jika dan hanya jika entri diagonalnya
berbeda berpasangan (lalu $x = (1, \dots, 1)$ berlaku, lewat
Vandermonde).
\end{solution}

\begin{solution}{exo:b3:modules:10}
Smith: $M = Q\operatorname{diag}(d_1,\dots,d_n)P$; lalu
\cref{exo:b3:modules:4} memberikan $[\Z^n : M\Z^n] = \prod\abs{d_i} =
\abs{\det M}$. Jika $\det M = \pm 1$: rumus adjugat
$M^{-1} = (\det M)^{-1}\,{}^{t}\!\operatorname{com}(M)$ berentri
bulat, sehingga $M \in GL_n(\Z)$; sebaliknya $MM^{-1} = I$
memberikan $\det M \cdot \det M^{-1} = 1$ di $\Z$, jadi $\det M = \pm1$.

Subgrup berindeks $m$ di $\Z^2$: subgrup $L$ semacam itu berrank
$2$ (karena indeksnya berhingga) dan mempunyai basis tunggal dalam \emph{bentuk
normal Hermite} $\bigl(\begin{smallmatrix} a & b\\ 0 &
d\end{smallmatrix}\bigr)$: di sini $d$ dicirikan oleh $L \cap (\{0\}
\times \Z) = \{0\} \times d\Z$, $a$ oleh $\pi_1(L) = a\Z$ (koordinat
pertamanya), dan $b$ lalu tunggal modulo $d$; normalkan $a, d
> 0$ dan $0 \leq b < d$. Indeksnya $ad = m$. Pencacahannya: untuk
setiap pembagi $d \mid m$ (dengan $a = m/d$), ada $d$ pilihan
$b$: totalnya $\sum_{d \mid m} d = \sigma_1(m)$.
\end{solution}

\begin{solution}{exo:b3:modules:11}
(a) Menurut teorema struktur, $G \cong \prod_i\Z/d_i\Z$, dan
$x^k = e$ terurai koordinat demi koordinat. Di $\Z/d\Z$: $kx \equiv 0
\pmod d$ mempunyai tepat $\gcd(k, d)$ penyelesaian (sebab $x$ haruslah
kelipatan $d/\gcd(k,d)$, dan ada $\gcd(k, d)$ buah di
antaranya). Lalu kalikan atas semua sukunya.

(b) Jika $G = \Z/d_s\Z$ siklik ($s = 1$), cacahannya
$\gcd(k, d_s) \leq k$. Jika $s \geq 2$: ambil $k = d_1$;
cacahannya $\prod_i\gcd(d_1, d_i) = d_1^{\,s} > d_1$ (sebab setiap
$\gcd$ sama dengan $d_1$ menurut rantai keterbagiannya): jadi persamaan
$x^{d_1} = e$ mempunyai lebih dari $d_1$ penyelesaian.

(c) Di dalam sebuah lapangan, $X^k - 1$ berakar paling banyak $k$
(\cref{ch:b3:rings}: polinomial tak nol berderajat $k$ atas
daerah integral), sehingga setiap subgrup berhingga $G \leq K^\times$
memenuhi kriteria (b): jadi $G$ siklik --- inilah bukti
struktural satu baris atas kesiklikan
$\mathbb F_q^\times$, yang melengkapi bukti pencacahan pada
\cref{ch:b3:galois}.
\end{solution}

\begin{solution}{exo:b3:modules:12}
(a) Jika $MN = I$ dengan $N$ bulat: $\det M\det N = 1$ dengan
keduanya bilangan bulat, sehingga $\det M = \pm1$. Sebaliknya jika $\det M =
\pm1$, rumus kofaktor $M^{-1} =
\frac1{\det M}\,{}^t\!\operatorname{com}(M)$ berentri bulat.

(b) Perkalian kiri dengan $E^{-k}$ mengurangkan $k$ kali baris
$2$ dari baris $1$; sedangkan dengan $F^{-k}$, $k$ kali baris $1$ dari baris
$2$. Diberikan $M = \bigl(\begin{smallmatrix}a & *\\ c &
*\end{smallmatrix}\bigr) \in SL_2(\Z)$, kolom pertamanya
$(a, c)$ merupakan vektor unimodular (sebab $\gcd(a, c) = 1$: ia
membagi $\det M = 1$). Jalankan algoritma Euclid pada $(a, c)$ lewat operasi
baris tersebut: setelah berhingga langkah kolomnya menjadi
$(\pm1, 0)$. Matriksnya kini berupa $\pm
\bigl(\begin{smallmatrix}1 & b\\ 0 & 1\end{smallmatrix}
\bigr) = \pm E^{b}$ (determinannya tetap $1$). Tersisa
menuliskan $-I$ dalam pembangunnya: $(EF^{-1}E)^2 =
\bigl(\begin{smallmatrix}0 & 1\\ -1 &
0\end{smallmatrix}\bigr)^2 = -I$ (periksalah kuadrat
matriks rotasi itu). Setelah diurai, $M$ merupakan kata dalam $E^{\pm1},
F^{\pm1}$.

(c) Algoritma Smith (\cref{met:b3:modules:smith}) memakai
persis operasi baris dan kolom semacam itu (ditambah penukaran dan pergantian
tanda, yang juga hasil kali operasi elementer sampai pada
tanda determinannya): jadi atas $\Z$, setiap matriks berbentuk
$U\,D\,V$ dengan $U, V$ hasil kali matriks elementer dan
$D$ bentuk Smithnya --- (b) adalah kasus $2\times2$ berdeterminan $1$
dari kenyataan umum bahwa matriks bertipe $E$ membangun
$SL_n(\Z)$.
\end{solution}

\begin{solution}{pb:b3:modules:1}
\textbf{1.} Pemetaan $\varphi$ aditif dan linear atas $K[X]$: $\varphi(X
\cdot (Q_i)_i) = \sum_i (XQ_i)(M)e_i = M\sum_i Q_i(M)e_i = X
\cdot \varphi\bigl((Q_i)_i\bigr)$, sebab struktur \omterm{def:b3:modules:module}{modul} $V_M$
adalah $X \cdot v = Mv$. Ia surjektif: vektor konstan sudah memberikan
seluruh $K^n$. Kolom ke-$j$ dari $XI - M$ adalah $Xe_j - \sum_i
m_{ij}e_i$, yang petanya $Me_j - \sum_i m_{ij}e_i = 0$.

\textbf{2.} Modulo kolom-kolomnya, $Xe_j \equiv \sum_i m_{ij}e_i$:
jadi sembarang vektor polinomial tereduksi, lewat induksi pada derajat
tertingginya, menjadi vektor \emph{konstan} $c \in K^n$. Jika vektor
semula terletak di $\ker\varphi$, maka $\varphi(c) = c = 0$ (pada
vektor konstan, $\varphi$ berupa pengenalan $K^n = V_M$), sehingga
vektor itu terletak di rentang kolomnya: $\ker\varphi = (XI_n -
M)K[X]^n$. Karena itu $V_M \cong K[X]^n/(XI - M)K[X]^n$, dan reduksi Smith
atas $K[X]$ (semua \omterm{thm:b3:modules:smith}{faktor invariannya} tak nol, sebab hasil kalinya
$\det(XI - M) = \chi_M$) memberikan $V_M \cong \bigoplus_i
K[X]/(f_i)$: yang $f_i$ tak konstan adalah \omterm{thm:b3:modules:frobenius}{invarian
keserupaannya}, yang dapat dihitung sebagai hasil bagi FPB minornya
(\cref{thm:b3:modules:smith}).

\textbf{3.} Untuk $\lambda I_n$: matriks $XI - \lambda I$ sudah berbentuk Smith:
invariannya $(X - \lambda, \dots, X - \lambda)$, sebanyak $n$ buah.
Entri diagonal yang berbeda: $V \cong \bigoplus_i K[X]/(X -
\lambda_i)$ dengan modulus yang \omterm{thm:b3:rings:crt}{komaksimal} berpasangan, sehingga sisa Tiongkok memampatkannya menjadi
satu \omterm{def:b3:modules:module}{modul} siklik $K[X]/\bigl(\prod_i(X - \lambda_i)\bigr)$:
satu invarian saja, yaitu $\chi$. Untuk $J_n(0)$: $\mu = X^n = \chi$ memaksa
satu invarian $X^n$. Untuk $\operatorname{diag}(J_2(0), J_1(0))$:
\omterm{thm:b3:modules:structure}{pembagi elementernya} $X^2, X$: jadi invariannya $P_1 = X \mid P_2 =
X^2$.

\textbf{4.} Pemetaan $P \mapsto P(u)$ memetakan $K[X]$ pada $K[u]$, dengan
kernel $(\mu_u)$ menurut definisi polinomial minimal: jadi $K[u]
\cong K[X]/(\mu_u)$, berdimensi $\deg\mu_u = \deg P_s = n_s$.

\textbf{5.} Kelinearan atas $K[X]$ memaksa $f(\bar Q) = f(Q\cdot\bar 1)
= Q\,f(\bar 1)$. Kelas $\bar 1$ memenuhi $P \bar 1 = 0$, sehingga
$P f(\bar 1) = 0$ menjadi keharusan. Sebaliknya jika $P\bar c = 0$,
maka $f(\bar Q) = Q\bar c$ terdefinisi dengan baik (sebab $Q \equiv Q' \bmod P
\Rightarrow (Q - Q')\bar c \in$ kelipatan $P\bar c = 0$) dan
linear atas $K[X]$.

\textbf{6.} Misalkan $g = \gcd(P, Q)$, $P = gP'$, $Q = gQ'$ dengan
$\gcd(P', Q') = 1$. Di $K[X]/(Q)$: $P\bar c = 0 \iff Q \mid Pc
\iff Q' \mid P'c \iff Q' \mid c$ (menurut Euclid, sebab $\gcd(P', Q') = 1$).
Jadi $\bar c$ yang diperbolehkan membentuk submodul yang dibangun oleh
$\bar{Q'}$, yang anihilatornya $\{R : Q \mid RQ'\} = (g)$:
sehingga submodul itu $\cong K[X]/(g)$. Bersama pertanyaan 5,
$\operatorname{Hom}(K[X]/(P), K[X]/(Q)) \cong K[X]/(\gcd(P,Q))$,
yang berdimensi $\deg\gcd(P,Q)$.

\textbf{7.} Pemetaan $v$ komutatif dengan $u$ jika dan hanya jika $v$ komutatif dengan setiap
$P(u)$, yaitu jika dan hanya jika $v$ linear atas $K[X]$: jadi $\mathcal C(u) =
\operatorname{End}_{K[X]}(V)$. Dengan menuliskan morfisma $V =
\bigoplus_j V_j$ sebagai matriks $(f_{ij})$, $f_{ij} \in
\operatorname{Hom}(V_j, V_i)$ (komposisikan dengan injeksi dan
proyeksinya), pertanyaan 6 memberikan
\[
\dim \mathcal C(u) = \sum_{i,j} \deg\gcd(P_i, P_j)
= \sum_{i,j} n_{\min(i,j)}
= \sum_{k=1}^{s} \bigl(2(s - k) + 1\bigr)\,n_k ,
\]
dengan memakai rantai keterbagiannya ($\gcd(P_i, P_j) =
P_{\min(i,j)}$) dan, untuk langkah terakhirnya, bahwa $\min(i,j) = k$
terjadi tepat pada $2(s - k) + 1$ pasangan $(i, j)$.

\textbf{8.} Karena $2(s-k)+1 \geq 1$, berlaku $\dim\mathcal C(u) \geq
\sum_k n_k = n$, dengan kesamaan jika dan hanya jika $s = 1$, yakni jika dan hanya jika $u$
siklik (\cref{exo:b3:modules:9}). Untuk $u = \lambda\,\mathrm{id}$:
$s = n$ dan semua $n_k = 1$: $\dim = \sum_{k=1}^n (2(n-k)+1) = n^2$
--- benar, sebab $\mathcal C(\lambda\,\mathrm{id}) = \mathcal
L(V)$.

\textbf{9.} Lewat rumusnya: invariannya $(X, X^2)$, jadi $s = 2$, $n_1 =
1$, $n_2 = 2$: $\dim = 3\cdot1 + 1\cdot2 = 5$. Secara langsung: dalam
basis $(e_1, e_2, e_3)$ dengan $u e_2 = e_1$, $ue_1 = ue_3 = 0$,
menuliskan $Au = uA$ untuk $A = (a_{ij})$ melahirkan syarat
$a_{21} = a_{23} = a_{31} = 0$ dan $a_{11} = a_{22}$: jadi ada lima
parameter \omterm{def:b3:modules:free}{bebas} $a_{11}{=}a_{22}, a_{12}, a_{13}, a_{32}, a_{33}$.

\textbf{10.} Polinomial $P(u)$ komutatif dengan apa pun yang
komutatif dengan $u$ (sebab ia jumlah pangkat $u$): jadi $K[u]
\subseteq \mathcal C(\mathcal C(u))$. Dan $u \in \mathcal C(u)$,
sehingga setiap $w \in \mathcal C(\mathcal C(u))$ komutatif dengan $u$.

\textbf{11.} Misalkan $V = K[u]x$ dan $v \in \mathcal C(u)$. Tulis
$v(x) = P(u)x$ (karena kesiklikan). Untuk $y = Q(u)x$ sembarang: $v(y) =
vQ(u)x = Q(u)v(x) = Q(u)P(u)x = P(u)y$. Jadi $v = P(u)$: $\mathcal
C(u) = K[u]$. Lalu $\mathcal C(\mathcal C(u)) = \mathcal
C(K[u]) = \mathcal C(u) = K[u]$ (komutatif dengan seluruh $K[u]$ sama saja
dengan komutatif dengan $u$). Jadi teoremanya berlaku pada kasus
siklik.

\textbf{12.} Proyeksi $\pi_i$ linear atas $K[X]$ (sebab penguraiannya
jumlah langsung submodul), jadi $\pi_i \in \mathcal C(u)$, dan $w$
komutatif dengannya: $w(V_i) = w\pi_i(V) = \pi_i w(V) \subseteq
V_i$. Pembatasan $w_i = w\restriction_{V_i}$ komutatif dengan
$u_i = u\restriction_{V_i}$ yang siklik (pertanyaan 10), dan $w_i
\in \mathcal C(u_i) = K[u_i]$ (pertanyaan 11): jadi $w_i = Q_i(u_i) =
Q_i(u)\restriction_{V_i}$ untuk suatu $Q_i \in K[X]$.

\textbf{13.} Keterdefinisian $\eta_{ij}(R(u)x_j) = R(u)x_i$:
jika $R(u)x_j = 0$ maka $P_j \mid R$, dan $P_i \mid P_j$ memberikan
$P_i \mid R$, sehingga $R(u)x_i = 0$ (sebab $\operatorname{Ann}(x_i) =
(P_i)$). Lalu $\eta_{ij}$ linear atas $K[X]$ menurut konstruksinya, dan
$\tilde\eta_{ij} = \eta_{ij}\pi_j$ merupakan komposisi
morfisma-$K[X]$ $V \to V$: jadi $\tilde\eta_{ij} \in \mathcal
C(u)$.

\textbf{14.} Nilaikan $w\tilde\eta_{ij} = \tilde\eta_{ij}w$ pada $x_j$:
ruas kirinya $w(x_i) = Q_i(u)x_i$; ruas kanannya
$\eta_{ij}\bigl(Q_j(u)x_j\bigr) = Q_j(u)x_i$. Karena itu $(Q_i -
Q_j)(u)\,x_i = 0$: jadi $P_i \mid Q_i - Q_j$ untuk setiap $i \leq j$.
Khususnya, dengan $Q = Q_s$: $Q \equiv Q_i \pmod{P_i}$, sehingga $Q(u)$
dan $Q_i(u)$ bersesuaian pada $V_i$ (yang dimatikan $P_i(u)$). Karena itu $w
= Q(u)$ pada setiap $V_i$, dan karenanya pada $V$: $\mathcal C(\mathcal C(u))
\subseteq K[u]$, lalu bersama pertanyaan 10, $\mathcal C(\mathcal
C(u)) = K[u]$.

\textbf{15.} Pernyataan pertamanya adalah pertanyaan 10--14 (atau, untuk
$u$ siklik, pertanyaan 11 saja). Untuk $u = \mathrm{id}$ dengan $n \geq
2$: $\mathcal C(u) = \mathcal L(V)$ berdimensi $n^2$, sedangkan
$\dim K[u] = \deg\mu_u = 1$. Jadi $\mathcal C(u) = K[u]$ gagal
telak; namun teorema bikomutannya tetap
berlaku ($\mathcal C(\mathcal L(V)) = K\,\mathrm{id} = K[u]$: pusat
aljabar matriksnya adalah skalar). Kesiklikanlah yang
membuat komutan \emph{tunggalnya} sudah berupa polinomial; sedangkan
komutan \emph{gandanya} selalu berupa polinomial.

\textbf{16.} Jika $P(XI - M)Q = S$ merupakan reduksi Smith (dengan $P, Q$
terbalikkan atas $K[X]$), mentransposkannya memberikan ${}^tQ\,(XI -
{}^tM)\,{}^tP = {}^tS = S$: bentuk Smith yang sama, sehingga $XI - M$ dan
$XI - {}^tM$ mempunyai \omterm{thm:b3:modules:smith}{faktor invarian} yang sama, yakni $M$ dan
${}^tM$ mempunyai \omterm{thm:b3:modules:frobenius}{invarian keserupaan} yang sama
(\cref{cor:b3:modules:descent}): jadi keduanya serupa.

\textbf{17.} \omterm{thm:b3:modules:frobenius}{Invarian keserupaan} $M$ atas $L$ adalah
\omterm{thm:b3:modules:smith}{faktor invarian} $XI - M$ di $L[X]$. Sebuah reduksi Smith atas
$XI - M$ di $K[X]$ --- dengan $P, Q$ terbalikkan atas $K[X]$,
diagonal, dan memenuhi rantai keterbagian --- \emph{juga} merupakan reduksi
Smith yang sah atas $L[X]$ (sebab $P, Q$ tetap terbalikkan: determinannya
konstanta tak nol), dan \omterm{thm:b3:modules:smith}{faktor invarian} yang monik bersifat tunggal: jadi
\omterm{thm:b3:modules:smith}{faktor invarian} yang dihitung atas $K$ dan atas $L$ berimpit. Karena itu $M \sim_L N$ jika dan hanya jika keduanya mempunyai \omterm{thm:b3:modules:smith}{faktor
invarian} yang sama, yaitu jika dan hanya jika $M \sim_K N$. Khususnya
matriks real yang sekawan atas $\C$ juga sekawan atas $\R$ --- pernyataan
yang kerap dibuktikan secara analitik (dengan mengkhususkan sebuah $P + \iu
Q$ yang terbalikkan), di sini dibuktikan secara struktural.

\textbf{18.} Uraikan $u = \bigoplus J_{m_t}(0)$ atas
blok Jordan nilpoten. Pada satu blok berukuran $m$ berlaku
$\operatorname{rk}J_m^k = \max(m - k, 0)$, sehingga
$\operatorname{rk}J_m^{k-1} - \operatorname{rk}J_m^k = 1$ bila $m
\geq k$, dan $0$ bila tidak. Dengan menjumlahkan atas bloknya:
$\operatorname{rk}u^{k-1} - \operatorname{rk}u^k =
\#\{t : m_t \geq k\}$. Karena itu barisan ranknya menentukan
multihimpunan $(m_t)$ --- yaitu sebuah partisi $n$ --- dan sebaliknya
setiap partisi terwujudkan: jadi kelas nilpoten $\leftrightarrow$
partisi $n$, atas lapangan apa pun. Untuk $M_5$: ada $p(5) = 7$ kelas
($5$; $4{+}1$; $3{+}2$; $3{+}1{+}1$; $2{+}2{+}1$;
$2{+}1{+}1{+}1$; $1^5$).

\textbf{19.} (a) Atas $\Q$ (atau lapangan berkarakteristik $0$ mana pun),
$D^n = 0$ dan $D^{n-1}(X^{n-1}) = (n-1)!\, \neq 0$: jadi $D$
nilpoten berindeks $n$ pada ruang berdimensi $n$, sehingga
siklik dengan invarian tunggal $X^n$ (sebab $x = X^{n-1}$ membangunnya:
turunan berulangnya merentang seluruh ruang). (b) Atas $\mathbb F_p$ dengan $p
< n$: $D^p = 0$, sebab turunan ke-$p$ setiap
monomial $X^m$ membawa faktor $m(m-1)\cdots(m-p+1)$, yaitu
hasil kali $p$ bilangan bulat berurutan, sehingga $\equiv 0 \bmod p$.
Tulis $n = ap + r$, $0 \leq r < p$. Maka $\ker D^k$ direntang
oleh monomial $X^m$ yang memenuhi $D^kX^m = 0$; dengan mencacah eksponen
$m < n$ menurut sisanya modulo $p$: $\dim\ker D^k = ak + \min(r,
k)$ untuk $0 \leq k \leq p$, sehingga $\operatorname{rk}D^{k-1} -
\operatorname{rk}D^k = \dim\ker D^k - \dim\ker D^{k-1} = a +
\mathbf 1_{k \leq r}$. Menurut pertanyaan 18, partisinya mempunyai $a$
blok berukuran tepat $p$ dan (bila $r > 0$) satu blok berukuran
$r$: jadi \omterm{thm:b3:modules:frobenius}{invarian keserupaannya} $X^r \mid X^p \mid \dots \mid X^p$.
Karakteristik mengubah bentuk kanonik operator yang paling akrab
dalam matematika.

\textbf{20.} Matriks $M \in GL_2$ mempunyai $s \in \{1, 2\}$ \omterm{thm:b3:modules:smith}{faktor
invarian}. Jika $s = 2$: $P_1 = P_2 = X - a$ (dengan $a \neq 0$ karena
keterbalikan), yakni $M = aI$, jadi berada di pusat. Jika $s = 1$: $M$
siklik dengan polinomial karakteristik $=$ polinomial minimal $\chi$
berderajat $2$, dan kelasnya bersesuaian dengan $\chi$ yang mungkin
dengan $\chi(0) \neq 0$: $\chi$ terpecah dengan akar berbeda $a \neq
b$ (pendampingnya $\sim$ diagonal); $\chi = (X - a)^2$ (pendamping,
tak semisederhana); dan $\chi$ \omterm{def:b3:rings:divisibility}{tak tereduksi}. Tepat satu jenis untuk masing-masing ---
sebab \omterm{thm:b3:modules:smith}{faktor invariannya} merupakan invarian yang lengkap.

\textbf{21.} Di pusat: ada $q - 1$ pilihan $a$. Nilai eigen terpecah
yang berbeda: pasangan tak terurut $\{a, b\} \subseteq
\mathbb F_q^\times$ dengan $a \neq b$: ada $\binom{q-1}2$ kelas.
Berpolinomial minimal $(X-a)^2$: ada $q - 1$ kelas. Polinomial kuadrat \omterm{def:b3:rings:divisibility}{tak tereduksi} dengan
suku konstan tak nol: semua polinomial kuadrat \omterm{def:b3:rings:divisibility}{tak tereduksi} memenuhinya
(sebab akarnya tak nol), dan ada $\frac{q^2 - q}2$
polinomial kuadrat monik yang \omterm{def:b3:rings:divisibility}{tak tereduksi} ($q^2$ polinomial kuadrat monik dikurangi
$\binom q2 + q = \frac{q^2+q}2$ yang terpecah). Totalnya:
\[
(q - 1) + \frac{(q-1)(q-2)}2 + (q - 1) + \frac{q^2 - q}2
= q^2 - 1 .
\]

\textbf{22.} Jika $M$ tidak siklik, maka $s = 2$ dan $M$ berupa skalar
(dikotomi pertanyaan 20 berlaku di $M_2$, baik terbalikkan maupun tidak: dua
\omterm{thm:b3:modules:smith}{faktor invarian} berderajat $1$ dengan $P_1 \mid P_2$ dan $P_1 =
P_2$ memaksa $M = aI$). Skalarnya berjumlah $q$ di antara $q^4$
matriks: jadi peluang kesiklikannya $1 - q^{-3}$. Secara umum
tempat kedudukan tak siklik di $M_n$ adalah tempat minor $(n-1)\times(n-1)$
dari $XI - M$ berbagi sebuah faktor --- sebuah syarat aljabar sejati
--- sehingga proporsinya sekecil $O(1/q)$ untuk $q$ besar: matriks
dengan $\chi = \mu$ adalah aturannya, dan argumen perekatan pada Bagian III
merupakan harga yang dibayar untuk perkecualiannya.

\textbf{23.} Unsur pusat $\mathcal C(u)$ terletak
di $\mathcal C(u)$ dan komutatif dengan setiap unsur
$\mathcal C(u)$, yakni terletak di $\mathcal C(\mathcal C(u)) =
K[u]$ (pertanyaan 14). Sebaliknya $K[u] \subseteq \mathcal C(u)$,
dan setiap $P(u)$ komutatif dengan setiap $v \in \mathcal C(u)$
(sebab $v$ semacam itu komutatif dengan $u$, sehingga dengan setiap pangkat $u$):
jadi $K[u]$ berada di pusat $\mathcal C(u)$. Karena itu $Z(\mathcal C(u))
= K[u]$. Akibatnya $\mathcal C(u)$ komutatif jika dan hanya jika
$\mathcal C(u) = Z(\mathcal C(u)) = K[u]$; dalam hal itu
$\dim\mathcal C(u) = \dim K[u] = n_s \leq n$, sedangkan pertanyaan 8
memberikan $\dim\mathcal C(u) \geq n$: sehingga $n_s = n$ dan $s = 1$,
yakni $u$ siklik. Sebaliknya, untuk $u$ siklik pertanyaan 11
memberikan $\mathcal C(u) = K[u]$, yang komutatif. Secara struktural: sebuah
aljabar matriks yang sama dengan pusatnya sendiri tidak lain adalah
aljabar polinomial $K[u]$ dari sebuah $u$ yang siklik.

\textbf{24.} Setiap koefisien $2s - 2i + 1$ pada rumus
Frobenius bernilai ganjil, sehingga
\[
\dim\mathcal C(u) = \sum_{i=1}^s(2s - 2i + 1)\,n_i
\equiv \sum_{i=1}^s n_i = n \pmod 2 .
\]
Untuk $n = 4$, barisan derajat $n_1 \leq \dots
\leq n_s$ yang mungkin bagi \omterm{thm:b3:modules:smith}{faktor invariannya}, dengan jumlah $4$, adalah $(4)$,
$(1, 3)$, $(2, 2)$, $(1, 1, 2)$, $(1, 1, 1, 1)$; masing-masing
terwujudkan, misalnya oleh matriks nilpoten dengan $P_i = X^{n_i}$ (rantai
keterbagiannya berlaku otomatis). Rumus itu berturut-turut memberikan
\[
1\cdot4 = 4, \quad 3 + 3 = 6, \quad 6 + 2 = 8, \quad
5 + 3 + 2 = 10, \quad 7 + 5 + 3 + 1 = 16 .
\]
Jadi himpunan nilainya $\{4, 6, 8, 10, 16\}$: bilangan genap dengan
keparitasan yang benar, tetapi $12$ dan $14$ tak pernah muncul --- antara
barisan yang hampir siklik dan $n^2$ milik skalar terdapat sebuah jurang.

\textbf{25.} $\abs{GL_2(\mathbb F_3)} = (q^2 - 1)(q^2 - q) =
8 \cdot 6 = 48$. Jenis pusat: $I$ dan $2I$, dua kelas
berukuran $1$. Pada ketiga jenis lainnya $M$ bersifat siklik (pertanyaan
20), sehingga \omterm{ex:b3:groups:actions}{sentralisatornya} di $GL_2$ adalah grup unsur terbalikkan
dari $\mathcal C(M) = K[M]$ (pertanyaan 11), dan ukuran
kelasnya $=48/\abs{K[M]^\times}$ menurut orbit--stabilisator. Nilai eigen
terpecah yang berbeda: hanya pasangan $\{1, 2\}$, jadi satu kelas; di sini $K[M]
\cong \mathbb F_3 \times \mathbb F_3$ (sisa Tiongkok pada $\chi = (X-1)
(X-2)$), unitnya $2 \cdot 2 = 4$, ukurannya $48/4 = 12$. Berpolinomial minimal
$(X - a)^2$ dengan $a \in \{1, 2\}$: dua kelas; di sini $K[M] \cong
\mathbb F_3[X]/((X-a)^2)$, unitnya $q^2 - q = 6$ (sebab suku konstan
unitnya $\neq 0$ setelah dipusatkan), ukurannya $48/6 = 8$.
Untuk $\chi$ \omterm{def:b3:rings:divisibility}{tak tereduksi}: polinomial kuadrat monik \omterm{def:b3:rings:divisibility}{tak tereduksi} atas
$\mathbb F_3$ berjumlah $(9 - 3)/2 = 3$, yakni
\[
X^2 + 1, \qquad X^2 + X + 2, \qquad X^2 + 2X + 2,
\]
(tanpa akar di $\mathbb F_3$: periksalah $0, 1, 2$); jadi tiga kelas dengan
$K[M] \cong \mathbb F_9$, unitnya $q^2 - 1 = 8$, ukurannya $48/8 =
6$. Persamaan kelasnya: $2\cdot1 + 1\cdot12 + 2\cdot8 + 3\cdot6 =
2 + 12 + 16 + 18 = 48$; dan ada $2 + 1 + 2 + 3 = 8 = q^2 - 1$
kelas, cocok dengan pertanyaan 21.

\end{solution}
